% =====================================================================
\section{Validación de requisitos}
\label{sec:validacion}

\subsection{Inspección formal de la Unidad IV}

\begin{table}[H]
\centering\footnotesize
\caption{Parámetros y resultados de la inspección formal tipo Fagan}
\label{tab:inspeccion}
\begin{tabular}{|L{6.4cm}|L{8.6cm}|}
\hline
\rowcolor{grisclaro}\textbf{Parámetro} & \textbf{Valor}\\ \hline
Documento inspeccionado & ERS/SRS del SIMPA v1.0 (Entrega 2A)\\ \hline
Alcance & \S1 a \S7; apéndices A--E excluidos\\ \hline
Páginas efectivas & 63 (pp. 8--70)\\ \hline
Roles & Moderador, Lector, Inspector 1 e Inspector 2, con acumulación declarada de dos roles en un integrante\\ \hline
Conformación del equipo en la inspección & Tres integrantes (conformación vigente en la Unidad IV)\\ \hline
Defectos únicos consolidados & 33\\ \hline
Distribución por severidad & 4 críticos, 24 mayores, 5 menores\\ \hline
Distribución por tipo & Consistencia 11, completitud 6, verificabilidad 5, ambigüedad 4, trazabilidad 4, factibilidad 3\\ \hline
Densidad de defectos & 0,52 defectos por página\\ \hline
Esfuerzo & 14,5 horas-persona; 0,44 horas-persona por defecto detectado\\ \hline
\end{tabular}
\end{table}

\noindent
Una precisión sobre la conformación del equipo antes de leer los resultados. La inspección se ejecutó
con los tres integrantes que componían el equipo de práctica experimental en la Unidad IV, lo que
obligó a acumular dos de los cuatro roles exigidos en un mismo participante; la acumulación se
declaró en su momento y no se disimula aquí. Con posterioridad a la inspección, el docente reagrupó
los equipos de práctica experimental conforme a los equipos del Proyecto Fin de Curso, con lo que
este equipo pasó a seis integrantes y la acumulación dejó de ser necesaria para trabajos ulteriores.
Los parámetros de la tabla anterior, la tasa de detección por rol y las métricas de inspección
corresponden a la ejecución original y no se recalculan: miden un hecho ocurrido en una fecha
determinada y con unos participantes determinados. Recalcularlos sobre la conformación actual
supondría atribuir a seis personas una detección que realizaron tres.

\vspace{0.3cm}
\noindent
Dos lecturas de esos números importan más que los números mismos.

La primera es la distribución por tipo. Once de los treinta y tres defectos son de consistencia y
ninguno es de sintaxis o de redacción. El documento sabía qué debía especificar; lo que fallaba era
que las mismas cifras se declaraban de forma distinta en secciones distintas. Es la firma de un
documento crecido por adición, sin una pasada de reconciliación final, y no la de un documento mal
escrito.

La segunda es el esfuerzo por defecto. Detectar un defecto costó menos de media hora-persona, cifra
que sostiene el argumento de Fagan de que la inspección es más barata que la corrección tardía
\cite{fagan1976}. El defecto \id{DEF-02} ---la contradicción entre 21 y 19 requisitos \emph{Must}---
habría llegado al diseño y habría producido dos planificaciones incompatibles.

\subsection{Corrección y gestión del cambio}

De los veintiocho defectos de severidad crítica y mayor, veinticinco se corrigieron directamente
sobre las fuentes y tres se tramitaron ante el Change Control Board por exceder el alcance de una
corrección editorial: su reparación consistía en añadir alcance al sistema, no en enmendar un texto.

\begin{table}[H]
\centering\footnotesize
\caption{Defectos derivados al Change Control Board}
\label{tab:def-ccb}
\begin{tabular}{|C{1.6cm}|C{1.6cm}|L{5.4cm}|L{6.0cm}|}
\hline
\rowcolor{grisclaro}\textbf{Defecto} & \textbf{RFC} & \textbf{Naturaleza y contenido} & \textbf{Decisión}\\ \hline
\id{DEF-13} & \id{RFC-01} & Alcance: \id{RF-36} y \id{RF-37} son \emph{Must} y carecían de historia de usuario & Aprobado por unanimidad\\ \hline
\id{DEF-16} & \id{RFC-02} & Calidad: la cobertura de las nueve características de ISO/IEC 25010:2023 no se sostenía & Aprobado por unanimidad\\ \hline
\id{DEF-12} & \id{RFC-03} & Normativa: \id{RL-04} a \id{RL-06} enunciaban obligaciones de la LOPDP sin requisito que las implementara & Aprobado con condiciones\\ \hline
\end{tabular}
\end{table}

\noindent
La deliberación sobre \id{RFC-03} merece registrarse porque el Change Control Board aprobó un cambio
que \emph{empeora} una métrica del propio proyecto. Incorporar \id{RF-40} a \id{RF-42} eleva los
requisitos \emph{Must} de veintiuno a veinticuatro y reduce la cobertura declarada del prototipo del
38,1\,\% al 33,3\,\%. El argumento que decidió la votación fue que el cumplimiento de una obligación
legal no se negocia contra un indicador de avance, y que el riesgo de omitirla recae sobre la
organización. Un órgano de control que solo aprobara lo que mejora sus indicadores no estaría
gobernando alcance sino administrando apariencia.

\subsection{Registro completo de defectos}

\begin{longtable}{|M{1.5cm}|L{2.4cm}|L{7.4cm}|M{1.9cm}|M{1.5cm}|}
\caption{Registro consolidado de los 33 defectos detectados en la inspección formal}
\label{tab:defectos}\\
\hline
\rowcolor{grisclaro}\textbf{ID} & \textbf{Sección} & \textbf{Descripción abreviada} & \textbf{Tipo} & \textbf{Sev.}\\ \hline
\endfirsthead
\hline
\rowcolor{grisclaro}\textbf{ID} & \textbf{Sección} & \textbf{Descripción abreviada} & \textbf{Tipo} & \textbf{Sev.}\\ \hline
\endhead
\id{DEF-01} & \S 3 (intro) & El encabezado de la \S 3 anuncia "9 restricciones de diseño"; la Tabla 64 y la \S 7.1 declaran 10 (RD-01 a RD-10). & Consistencia & Mayor\\ \hline
\id{DEF-02} & \S 5.1 vs \S 5.3.1 & La \S 5.1 declara una distribución MoSCoW de 21 Must, 16 Should y 2 Could (39 RF); la \S 5.3.1 declara 19 Must, 14 Should y 2 Could (35 RF). Las dos cifras no... & Consistencia & Crítico\\ \hline
\id{DEF-03} & \S 5.3.3, Anexo C & Ambos remiten a un archivo con "la tabla completa para los 35 requisitos funcionales", cuando el ERS especifica 39. & Consistencia & Mayor\\ \hline
\id{DEF-04} & \S 5.4, Tabla 69, \S 7.1, Anexo D & El número de filas de la matriz de trazabilidad se declara alternativamente como 52 y como 48 dentro del mismo documento; el diccionario define TR-01 a TR-48. & Consistencia & Mayor\\ \hline
\id{DEF-05} & \S 2.2 vs \S 3.2 & La \S 2.2 agrupa las funciones del producto en seis bloques; la \S 3.2 desarrolla siete (Bloque I a Bloque VII). & Consistencia & Mayor\\ \hline
\id{DEF-06} & Tabla 52 (RD-01) & RD-01 se justifica con "el personal ya dispone de teléfono inteligente propio (EV-05, EV-06, EV-07)". La \S 2.7 declara esa suposición invalidada por EV-12... & Consistencia & Crítico\\ \hline
\id{DEF-07} & Tabla 52 (RD-02) & RD-02 se justifica con "ninguna persona respondiente declaró señal permanente (EV-10)"; la \S 2.6 y la Tabla 7 reportan que el 25,8\% sí declara señal siempre... & Consistencia & Mayor\\ \hline
\id{DEF-08} & Tabla 78 vs Tabla 24 & El catálogo de evidencias atribuye a EV-13 el origen de RF-29, pero la ficha de RF-29 declara como origen EV-06 y EV-07 únicamente. & Consistencia & Mayor\\ \hline
\id{DEF-09} & \S 5.2, \S 5.4.1 (Tabla 70) & Se afirma que los 19 requisitos Must tienen historia de usuario y caso de uso (100\%). Si los Must son 21, la cobertura declarada es falsa y quedan requisitos... & Consistencia & Crítico\\ \hline
\id{DEF-10} & Tabla 68 & La tabla dice presentar los diez requisitos ordenados de forma descendente por WSJF, pero omite RF-05 y RF-30, a los que la Tabla 71 asigna WSJF 4,80, superior... & Consistencia & Mayor\\ \hline
\id{DEF-11} & \S 5.3.3 & La lectura del resultado afirma que RF-35 y RF-10 encabezan el orden con WSJF 7,00, pero RF-10 no figura en la Tabla 68; su valor solo aparece en la Tabla 71. & Consistencia & Mayor\\ \hline
\id{DEF-12} & Tabla 51 (\S 3.4) & RL-04, RL-05 y RL-06 enuncian obligaciones de la LOPDP (exportación de datos personales, rectificación con bitácora, supresión al término de la relación... & Completitud & Mayor\\ \hline
\id{DEF-13} & \S 5.2 & RF-36 y RF-37, clasificados Must en sus fichas, no tienen historia de usuario ni criterio de aceptación entre HU-01 y HU-19, pese a que la \S 5.2 declara cubrir... & Completitud & Crítico\\ \hline
\id{DEF-14} & \S 1.2.3 & La lista de funcionalidades principales omite el ciclo de liquidación semanal (RF-36 a RF-39), que la \S 7.2 califica como proceso central de la organización. & Completitud & Menor\\ \hline
\id{DEF-15} & Tabla 73 (L-06) & L-06 declara que las sesiones de walkthrough son inferiores a las tres exigidas y que "se declara el número efectivo"; ese número no se declara en ninguna... & Completitud & Mayor\\ \hline
\id{DEF-16} & \S 3.3 & Se afirma que los RNF cubren las nueve características de ISO/IEC 25010:2023, pero se emplea "Portabilidad" como característica de primer nivel (en la edición... & Completitud & Mayor\\ \hline
\id{DEF-17} & Tabla 39 (RF-12) & El requisito incluye "condiciones anómalas" como disparador de alerta sin definir qué constituye una condición anómala; el criterio de verificación solo cubre... & Verificabilidad & Mayor\\ \hline
\id{DEF-18} & Tabla 34 (RF-16) & No se define la fórmula, las variables ni la escala del "indicador de desempeño". El criterio de verificación ("el sistema calcula un indicador") es... & Verificabilidad & Mayor\\ \hline
\id{DEF-19} & Tabla 45 (RF-18) & La estimación de producción carece de tolerancia, margen de error o valor de referencia, por lo que ningún resultado puede declararse incorrecto. & Verificabilidad & Mayor\\ \hline
\id{DEF-20} & Tabla 49 (RNF-10) & Una disponibilidad mensual $\geq$99\% equivale a un máximo de 7 h 12 min de indisponibilidad (720 h $\times$ 1\%), no a las 7 h 18 min declaradas. & Verificabilidad & Menor\\ \hline
\id{DEF-21} & Tabla 49 (RNF-03) & Se fijan 50 sesiones concurrentes sin justificación; la población declarada es de 62 personas con conectividad mayoritariamente intermitente (74,2\% sin señal... & Verificabilidad & Menor\\ \hline
\id{DEF-22} & Tabla 35 (RF-07) & La precondición exige que "la imagen cumple los requisitos mínimos de calidad", pero esos requisitos no se especifican en ninguna sección del documento. & Ambigüedad & Mayor\\ \hline
\id{DEF-23} & Tabla 40 (RF-22) & La descripción exige alertar cuando el porcentaje "se aproxime al umbral", sin cuantificar la proximidad; el criterio de verificación usa "alcance o supere",... & Ambigüedad & Mayor\\ \hline
\id{DEF-24} & Tabla 43 (RF-25) / Tabla 52 (RD-08) & "El valor recomendado para la variedad" no está tabulado en ninguna parte. La única cifra aparece en la leyenda de la Figura 10 ("24 h híbrido / menor en... & Ambigüedad & Mayor\\ \hline
\id{DEF-25} & Tabla 37 (RF-09) & Se exige recomendar la acción de control "según la etapa del insecto", pero las etapas del ciclo no están enumeradas ni acotadas; el criterio de verificación... & Ambigüedad & Menor\\ \hline
\id{DEF-26} & Tabla 69 vs Tabla 57 & La matriz enlaza RF-16, RF-17 y RF-20 a CU-05, pero la especificación textual de CU-05 declara únicamente RF-19 y RF-26. La trazabilidad no cierra en sentido... & Trazabilidad & Mayor\\ \hline
\id{DEF-27} & Tabla 69 vs Tabla 9 & La matriz asocia RF-29 y RF-36 a RF-39 con el prototipo MU-02, mientras que la Tabla 9 declara que MU-02 soporta únicamente RF-12, RF-19 y RF-26. & Trazabilidad & Mayor\\ \hline
\id{DEF-28} & \S 1.4 y Referencias & Se declara que el documento cita 31 fuentes primarias; las entradas [27], [28], [30] y [31] figuran en la lista de referencias sin aparecer citadas en el texto. & Trazabilidad & Mayor\\ \hline
\id{DEF-29} & Tabla 78 vs portada & EV-13 se fecha entre el 03 y el 07/08/2026, es decir, con recolección posterior a la fecha de versión del propio documento (03/08/2026). & Trazabilidad & Mayor\\ \hline
\id{DEF-30} & Tabla 49 (RNF-01, RNF-02) & Los valores objetivo se verifican sobre un conjunto etiquetado de $\geq$200 imágenes que L-05 declara inexistente, y el documento no incorpora en el alcance de la... & Factibilidad & Mayor\\ \hline
\id{DEF-31} & \S 3.2.4 (RF-14) / Tabla 49 (RNF-05) & RF-14 exige derivar el conteo de flores planta por planta a partir de marcaciones GPS y RNF-05 admite una precisión de hasta 10 m. El ERS no declara la... & Completitud & Mayor\\ \hline
\id{DEF-32} & Tabla 73 & Los identificadores de limitación no siguen orden correlativo (L-01...L-07, L-10, L-11, L-08, L-09) y L-07 se declara "Resuelta" pero permanece listada como... & Factibilidad & Menor\\ \hline
\id{DEF-33} & \S 6.1.1 (Tabla 71) & La tabla se titula "cobertura sobre los requisitos funcionales Must" pero incluye RF-20 y RF-15, que sus fichas clasifican como Should. El denominador de 19 y... & Factibilidad & Mayor\\ \hline
\end{longtable}

\subsection{Métricas de la inspección e interpretación}

\begin{table}[H]
\centering\footnotesize
\caption{Métricas de la inspección formal y su lectura}
\label{tab:met-inspeccion}
\begin{tabular}{|L{3.8cm}|L{3.6cm}|L{7.6cm}|}
\hline
\rowcolor{grisclaro}\textbf{Métrica} & \textbf{Valor} & \textbf{Qué revela del ERS}\\ \hline
Densidad de defectos & 0,52 def./página & Con más de un defecto cada dos páginas, el documento no falla en un punto aislado sino de forma distribuida: apunta a un problema de integración entre secciones redactadas por separado\\ \hline
Distribución por tipo & Consistencia 11; completitud 6; verificabilidad 5; ambigüedad 4; trazabilidad 4; factibilidad 3 & La consistencia concentra un tercio: el documento sabe qué especificar, pero las mismas cifras se declaran distinto en secciones distintas\\ \hline
Distribución por severidad & 4 críticos; 24 mayores; 5 menores & El predominio de la severidad mayor sobre la menor indica que los defectos son de contenido especificado y no de forma\\ \hline
Tasa de detección por rol & Inspector 2: 11; Inspector 1: 8; Moderador: 7; Lector: 7 & Los cuatro roles superan el mínimo de seis y ninguno concentra más del 33\,\%: el reparto por secciones funcionó\\ \hline
Esfuerzo & 14,5 h-persona; 0,44 h-persona por defecto & Detectar un defecto costó menos de media hora-persona, cifra que justifica repetir la inspección antes de la siguiente entrega\\ \hline
\end{tabular}
\end{table}

\subsection{Re-inspección de la Unidad V}

La métrica de corrección de la auditoría requiere medir cuántos defectos sobrevivieron a la
inspección. Esa medición exige una re-inspección sobre el documento ya corregido, y su resultado es
el siguiente.

\begin{longtable}{|C{1.6cm}|L{4.0cm}|L{5.4cm}|C{1.5cm}|L{2.8cm}|}
\hline
\rowcolor{grisclaro}\textbf{ID} & \textbf{Ubicación} & \textbf{Defecto residual} & \textbf{Sev.} & \textbf{Estado}\\ \hline
\endhead
\id{RES-01} & Apéndice, diccionario de identificadores & Declara el rango de requisitos funcionales como \id{RF-01} a \id{RF-39}, cuando \id{RFC-03} lo elevó a \id{RF-42} & Mayor & Corregido\\ \hline
\id{RES-02} & Apéndice, diccionario de identificadores & Declara el rango de no funcionales como \id{RNF-01} a \id{RNF-18}, cuando \id{RFC-02} incorporó \id{RNF-19} & Mayor & Corregido\\ \hline
\id{RES-03} & Tabla de recuentos de la \S5.1 & Declara 39 requisitos funcionales y 18 no funcionales, cifras anteriores a las RFC & Mayor & Corregido\\ \hline
\id{RES-04} & Apéndice C & \id{RNF-19} no figura en la tabla de priorización: es el único requisito del catálogo sin prioridad MoSCoW & Mayor & Corregido\\ \hline
\id{RES-05} & \S4, casos de uso & Ocho de los dieciocho casos de uso identificados carecen de especificación textual & Mayor & Corregido\\ \hline
\id{RES-06} & \S5.2, criterios de aceptación & Solo 24 de 61 requisitos tienen criterio en formato Dado--Cuando--Entonces & Mayor & Corregido\\ \hline
\id{RES-07} & \S2.3 y fichas & Dos partes interesadas se contaban como clases de usuario sin ser actor de ningún requisito & Menor & Corregido\\ \hline
\caption{Defectos residuales detectados en la re-inspección de la Unidad V}
\label{tab:residuales}
\end{longtable}

\subsection{Por qué escaparon: análisis de causa}

Los siete defectos residuales no se distribuyen al azar, y su patrón dice algo sobre el
procedimiento de inspección más que sobre el documento.

\textbf{Cuatro se concentran en los apéndices y en los conteos globales} (\id{RES-01} a \id{RES-04}).
El alcance declarado de la inspección de la Unidad IV fue \S1 a \S7, con los apéndices excluidos de
forma expresa. La exclusión era razonable ---los apéndices son material de referencia--- pero dejó
fuera precisamente el lugar donde se acumulan los recuentos que las tres RFC modificaron. La causa no
es descuido: es que el alcance se fijó antes de saber qué iban a cambiar las RFC, y nadie lo revisó
después.

\textbf{Dos son de completitud del conjunto, no de corrección de sus elementos} (\id{RES-05},
\id{RES-06}). La lista de verificación empleada preguntaba si cada elemento revisado era correcto, no
si el conjunto estaba completo. Un inspector que recorre diez casos de uso bien especificados
responde ``sí cumple'' a todas las preguntas, y ninguna le obliga a contar cuántos deberían existir.

\textbf{Uno es un defecto de la propia auditoría} (\id{RES-07}). La primera medición de la cobertura
de actores contó ocho actores mezclando la tabla de clases de usuario con el mapa de partes
interesadas, y arrojó 75\,\% cuando el valor real era 100\,\%. Se corrigió el conteo y se añadió al
ERS una nota que declara la distinción, de modo que la métrica no vuelva a medirse mal. Es el único
de los siete que no era un defecto del ERS sino del instrumento que lo mide, y se documenta como tal
en lugar de retirarlo del registro.

\subsection{Acción correctiva de proceso}

Las siete correcciones están aplicadas, pero la conclusión operativa no es haberlas aplicado. Es que
la lista de verificación de la inspección debe modificarse en dos puntos para la siguiente entrega:
incluir los apéndices y los conteos globales dentro del alcance cuando el documento haya sufrido
cambios aprobados, e incorporar preguntas sobre la completitud del conjunto y no solo sobre la
corrección de sus elementos. Corregir siete defectos mejora este documento; cambiar la lista de
verificación mejora todos los siguientes.
